\documentclass[11pt,oneside]{article}
\usepackage[T1]{fontenc}
\usepackage[utf8]{inputenc}
\usepackage{lmodern}
\usepackage[margin=1.03in,headheight=14pt]{geometry}
\usepackage{amsmath,amssymb,amsthm,mathtools}
\usepackage{microtype,booktabs,array,longtable,enumitem}
\usepackage[dvipsnames]{xcolor}
\usepackage{hyperref}
\usepackage[nameinlink,noabbrev]{cleveref}
\usepackage{fancyhdr}
\hypersetup{colorlinks=true,linkcolor=MidnightBlue,citecolor=MidnightBlue,urlcolor=MidnightBlue,
 pdftitle={Hartogs Coherence, Finite Maps, and Residue Duality over the Surcomplex Numbers},
 pdfauthor={Research manuscript prepared with ChatGPT},pdfsubject={Support-controlled several-variable Hahn-coherent analysis}}
\pagestyle{fancy}\fancyhf{}
\fancyhead[L]{\small Hartogs coherence and finite surcomplex maps}
\fancyhead[R]{\small\thepage}
\renewcommand{\headrulewidth}{0.3pt}
\setlength{\parskip}{0.3em}
\setlength{\parindent}{1.2em}
\setlist[enumerate]{leftmargin=*,itemsep=0.25em,topsep=0.3em}
\numberwithin{equation}{section}
\newtheorem{theorem}{Theorem}[section]
\newtheorem{lemma}[theorem]{Lemma}
\newtheorem{proposition}[theorem]{Proposition}
\newtheorem{corollary}[theorem]{Corollary}
\theoremstyle{definition}
\newtheorem{definition}[theorem]{Definition}
\newtheorem{example}[theorem]{Example}
\newtheorem{question}[theorem]{Question}
\theoremstyle{remark}\newtheorem{remark}[theorem]{Remark}
\newcommand{\C}{\mathbb C}
\newcommand{\N}{\mathbb N}
\newcommand{\Z}{\mathbb Z}
\newcommand{\No}{\mathbf{No}}
\newcommand{\SC}{\mathbf{SC}}
\newcommand{\OO}{\mathcal O}
\newcommand{\HH}{\mathcal H_{\Gamma}}
\newcommand{\Hp}{\mathcal H^{+}_{\Gamma}}
\newcommand{\Ip}{\mathcal I_{\Gamma}}
\newcommand{\mm}{\mathfrak m}
\newcommand{\sh}{^{\sharp}}
\newcommand{\supp}{\operatorname{supp}}
\newcommand{\st}{\operatorname{st}}
\newcommand{\ord}{\operatorname{ord}}
\newcommand{\Tr}{\operatorname{Tr}}
\newcommand{\Hom}{\operatorname{Hom}}
\newcommand{\Mat}{\operatorname{Mat}}
\newcommand{\Res}{\operatorname{Res}}
\newcommand{\rank}{\operatorname{rank}}
\newcommand{\NF}{\operatorname{NF}}
\newcommand{\id}{\operatorname{id}}
\newcommand{\hoint}{\mathop{\oint}\nolimits^{\mathrm H}}
\newcommand{\boxset}{\mathcal B_{\boldsymbol m}}
\newcommand{\Span}{\operatorname{span}}
\DeclareMathOperator{\Spec}{Spec}
\newcolumntype{L}[1]{>{\raggedright\arraybackslash}p{#1}}
\begin{document}
\thispagestyle{plain}
\begin{center}
{\LARGE\bfseries Hartogs Coherence, Finite Maps,\par
and Residue Duality\par
over the Surcomplex Numbers\par}
\vspace{0.8em}
{\large Support-controlled several-variable extensions\par
of Hahn-coherent analysis\par}
\vspace{0.8em}
{September 21, 2026}
\end{center}
\begin{abstract}
The supplied manuscript on surcomplex analysis develops a one-variable theory of Hahn-coherent holomorphic functions and identifies parameter-dependent division and finite mappings as unfinished directions. We establish a several-variable extension with explicit support control. First, a Baire-category argument shows that pointwise well-ordered supports of holomorphic coefficient functions force a common well-ordered support. Combined with classical Hartogs separate holomorphy, it yields an automatic-coherence theorem: separate Hahn coherence in every surcomplex coordinate slice implies joint Hahn coherence, without assuming a common slice support.

Second, we prove a fixed-domain matrix division theorem and parameter-dependent Weierstrass preparation. Iterated matrix division treats genuinely coupled systems
\(F_i(x,w)=w_i^{m_i}+E_i(x,w)\), where every \(E_i\) has positive Hahn support. Their quotient algebra is finite free of rank \(\prod_i m_i\), with an explicit rectangular monomial basis. All division coefficients remain holomorphic on the original ordinary parameter domain, and their supports lie in the input support plus the positive monoid generated by the perturbations. After specialization to a surcomplex parameter, the algebra decomposes into the formal local intersection algebras of the actual infinitesimal roots. Thus the total intersection multiplicity is exactly \(\prod_i m_i\).

Finally, a coefficientwise multidimensional contour functional descends to this algebra and gives a perfect relative residue pairing, including at colliding roots. We derive collision-stable trace formulas and work out a two-variable splitting example, a nonpolynomial example, and a nondiscrete-support example. The results address a precise portion of the program in the supplied manuscript; they do not assert a solution to an independently documented classical open conjecture. The exact formulations are proposed research contributions, with novelty qualified by the literature comparison and the absence of independent proof review.
\end{abstract}
\vspace{0.3em}
\noindent\textbf{Keywords.} Surcomplex numbers; Hahn series; separate holomorphy; Weierstrass division; infinitesimal deformation; finite flat algebra; intersection multiplicity; residue duality.

\clearpage
\setcounter{tocdepth}{1}
\tableofcontents
\vspace{1em}
\noindent\textbf{Reading guide.}
Sections~\ref{sec:foundations}--\ref{sec:operators} provide the shared support machinery.
The automatic-coherence result in Section~\ref{sec:hartogs} is independent of the
finite-intersection construction. For the latter, matrix division in
Section~\ref{sec:division} leads to the finite-free algebra in
Section~\ref{sec:finite}, its actual root multiplicities in
Section~\ref{sec:fibers}, and residue duality in Section~\ref{sec:residue}.
The examples in Section~\ref{sec:examples} can be read alongside those theorems;
Section~\ref{sec:literature} separates the proposed contributions from their
classical and non-Archimedean predecessors.
\clearpage

\section{The question addressed and the scope of the results}\label{sec:scope}
\subsection{A concrete continuation of the supplied framework}
The starting point is the supplied manuscript \emph{Surcomplex Analysis: Infinitesimal Calculus, Hahn-Coherent Holomorphy, and Contour Theory}, distributed as \texttt{surcomplex\_analysis(3).tex} \cite{Source}. Its section entitled ``Further mathematical targets suggested by the results'' specifically asks for coherent division with holomorphic parameters and finite-mapping theorems with globally stated support control. The present article addresses that direction. Its definitions of surcomplex numbers, strong summability, standard-part thickenings, and Hahn coherence follow that manuscript.

A typical datum is
\[
 F(x,w)=\sum_{\gamma\in S}f_\gamma(x,w)t^\gamma,
 \qquad t^\gamma=\omega^{-\gamma},
\]
where \(S\) is a well-ordered set of surreal exponents and every \(f_\gamma\) is an ordinary holomorphic function on the \emph{same} domain. No bound on the ordinary sizes of the functions \(f_\gamma\) is imposed. Evaluation at finite surcomplex points uses Taylor expansion in infinitesimal displacements and strong Hahn summation. This is neither ordinary convergence in a real parameter \(t\) nor convergence of partial sums in the full fine topology on the surcomplex class.

The one-variable preparation and moving-pole residue theorems in \cite{Source} are not claimed as new here. The central extension is the passage from one equation to coupled analytic systems, while preserving a fixed ordinary domain and one explicitly controlled support. A second, logically independent extension removes a common-support assumption from separate coherence.

\subsection{The two principal assertions}
For orientation, we state informal versions before introducing the complete notation.

\paragraph{Automatic Hartogs coherence.}
Suppose a surcomplex-valued function on a product of thickened ordinary plane domains is Hahn-coherent in each variable whenever the other variables are fixed at arbitrary finite surcomplex values. Then it is jointly Hahn-coherent. The global well-ordered support is a consequence, not an additional hypothesis. The precise statement is \cref{thm:hartogs}.

\paragraph{Finite intersections and duality with support bounds.}
Let \(x\) denote ordinary holomorphic parameters, let \(w=(w_1,\ldots,w_n)\) be fiber variables on a fixed polydisk, and suppose
\begin{equation}\label{eq:intromodel}
 F_i(x,w)=w_i^{m_i}+E_i(x,w),\qquad m_i\geq1,
 \qquad \supp(E_i)\subset\Gamma_{>0}.
\end{equation}
The errors can depend on \emph{all} variables. If \(E\) is their common positive support and \(M=E^*\) its additive monoid, then:
\begin{enumerate}
\item the quotient by \((F_1,\ldots,F_n)\) has the basis
\(w^\alpha\), \(0\leq\alpha_i<m_i\), over the parameter coefficient ring;
\item a numerator of support \(S\) has normal-form and division coefficients supported in \(S+M\), on the original domain;
\item each surcomplex parameter fiber has total actual local intersection multiplicity \(N=\prod_i m_i\), entirely in the infinitesimal fiber monad;
\item a coefficientwise residue functional induces an invertible \(N\times N\) pairing matrix, even when individual roots collide.
\end{enumerate}
These assertions are proved in \cref{thm:finitefree,thm:fiberlength,thm:duality}. They are not consequences of applying the one-variable root-count theorem separately to each equation: the equations are coupled. The matrix division step controls that coupling.

\subsection{What ``new'' means in this article}
The literature search compared the supplied category with surreal analytic structures, Hahn fields, analytic structures on valued fields, and existing Hahn-meromorphic function theory. The precise automatic-support Hartogs theorem and the combined fixed-domain, arbitrary-support finite-intersection and duality package were not located in that search. This is a statement about the search, not a proof of priority. There may be equivalent results under different terminology, or stronger general theorems from which some assertions can be derived.

The article supplies proofs of its stated results rather than announcing a solution to a named open problem whose open status has not been independently established. Classical ingredients are identified explicitly. The manuscript has not been independently refereed or formally verified. The symbolic checks in the accompanying archive test examples, not the general theorems.

\subsection{Three boundaries that matter}
Our finite-flat statements concern explicitly constructed coefficient algebras and their actual surcomplex fibers. They are not assertions that the fine topology makes a projection proper. The leading systems treated in full detail are coordinate powers, or coordinate powers multiplied by ordinary zero-free units. Arbitrary isolated ordinary complete intersections are not silently included. Finally, we do not construct a universal continuation theory at infinite scales, a global theory of essential singularities, or a general Noetherian analytic geometry over the entire surcomplex class.

\section{Hahn-coherent functions in several variables}\label{sec:foundations}
\subsection{Set-sized workspaces and valuation}
We use a class set theory admitting the usual treatment of the surreal class \(\No\), and put \(\SC=\No[i]\). Every support and every recursion in this article is set-sized. A set-sized additive subgroup \(\Gamma\subseteq\No\) determines the Hahn field
\[
 K_\Gamma=\C((t^\Gamma))\subseteq\SC,
 \qquad t^\gamma=\omega^{-\gamma}.
\]
The embedding is given by Conway normal forms. For a nonzero Hahn series,
\[
 v\left(\sum_{\gamma\in S}a_\gamma t^\gamma\right)
 =\min\{\gamma:a_\gamma\ne0\}.
\]
Finite elements have nonnegative valuation and infinitesimals have positive valuation. Every finite \(z\in\SC\) is uniquely \(c+\varepsilon\), with \(c\in\C\) and \(\varepsilon\) infinitesimal. We write \(\st(z)=c\) and denote the infinitesimals by \(\mm\).

Every set of surcomplex constants used in a proof lies in a suitable workspace, which may be enlarged. When algebraic closure is needed, we take the divisible hull of the exponent group. The Hahn field with divisible value group and algebraically closed residue field is algebraically closed; see \cite[Corollary 4]{Poonen}. No proper-class polynomial factorization or proper-class sum is required.

For an ordinary open set \(U\subseteq\C^d\), define
\[
 U\sh=\{z\in(\C+\mm)^d:\st(z)\in U\}.
\]
The standard part is taken coordinatewise. The fiber monad of zero is \(\mm^n\).

\subsection{Strong summation and positive monoids}
\begin{definition}[Strong summability]
A set-indexed family of Hahn series is strongly summable when the union of its supports is well ordered and each exponent occurs in only finitely many members of the family. The sum is formed by finite addition at each exponent.
\end{definition}

We use the following support facts, often called the Hahn--Neumann lemma; their use in this context is also explained in \cite{Source,Poonen}. If \(S,T\) are well-ordered subsets of an ordered abelian group, then \(S+T\) is well ordered, and each exponent has only finitely many representations \(s+t\). If \(E\subseteq\Gamma_{>0}\) is well ordered, then
\[
 E^*=\{e_1+\cdots+e_k:k\geq0,\ e_j\in E\}
\]
is well ordered. Each exponent has only finitely many representations by finite nondecreasing words in \(E\), hence only finitely many by ordered words. We include the empty word, with sum zero.

In particular, if \(M=E^*\), then \(M+M=M\). For any well-ordered \(S\), the set \(S+M\) controls arbitrary finite iterations that add letters from \(E\), and at a fixed exponent only finitely many such iterations can contribute. This finiteness, rather than a convergence assertion about an iterative sequence, is the mechanism of the proofs below.

\subsection{Coefficient rings}
For a connected ordinary domain \(U\), write
\begin{align*}
 \HH(U)&=\left\{\sum_{\gamma\in S}f_\gamma t^\gamma:
 S\subseteq\Gamma\text{ well ordered},\ f_\gamma\in\OO(U)\right\},\\
 \Hp(U)&=\{F\in\HH(U):\supp(F)\subseteq\Gamma_{\geq0}\},\\
 \Ip(U)&=\{F\in\HH(U):\supp(F)\subseteq\Gamma_{>0}\}.
\end{align*}
Zero coefficient functions are omitted from the exact support. In particular \(\Ip(U)\) is an ideal of \(\Hp(U)\), and
\[
 \Hp(U)/\Ip(U)\cong\OO(U).
\]
It is not an ideal of the full ring \(\HH(U)\). A subscript \(0\) denotes reduction to the coefficient at exponent zero when the element has nonnegative support.

For finite vectors and matrices, support means the union of the supports of the entries. On a connected domain \(\OO(U)\) is an integral domain, so taking the least exponent of a nonzero Hahn-holomorphic datum behaves as expected under products. We also use the same notation for open sets with several connected components by treating components separately.

\begin{remark}[Why we do not invoke Noetherianity]\label{rem:nonnoetherian}
If \(\Gamma\) is nonzero and divisible, the positive ideal in \(\Hp(\{\mathrm{pt}\})\) is not finitely generated. Indeed, finitely many positive-valued generators have a least valuation \(\delta>0\), whereas \(t^{\delta/2}\) cannot lie in their ideal. Thus a Noetherian completion theorem cannot simply be assumed for this integral coefficient ring. Our arguments give explicit splittings and inverses instead. This observation does not say that the field \(K_\Gamma\) is non-Noetherian.
\end{remark}

\subsection{Evaluation, identity, and specialization}
For \(F=\sum f_\gamma t^\gamma\in\HH(U)\), define
\begin{equation}\label{eq:multieval}
 F\sh(c+\varepsilon)
 =\sum_{\gamma\in S}\sum_{\alpha\in\N^d}
 \frac{\partial^\alpha f_\gamma(c)}{\alpha!}
 t^\gamma\varepsilon^\alpha.
\end{equation}
Here \(\N\) includes zero, and multi-index notation has its usual meaning.

\begin{lemma}[Evaluation and fixed-domain specialization]\label{lem:eval}
The family in \eqref{eq:multieval} is strongly summable. Evaluation is an injective algebra homomorphism. It commutes with partial differentiation in the ordinary variables, leaving the monomials \(t^\gamma\) constant.

If \(F\in\HH(U\times V)\) and \(a=c+\varepsilon\in U\sh\), evaluation only in the \(U\)-variables gives a Hahn-coherent datum on the same domain \(V\), in a workspace containing the displacement exponents. Its support is contained in
\[
 \supp(F)+\left(\bigcup_j\supp(\varepsilon_j)\right)^*.
\]
\end{lemma}
\begin{proof}
Set \(E=\bigcup_j\supp(\varepsilon_j)\). This is a finite union of well-ordered positive sets. Every term in the evaluation has support in \(S+E^*\). A fixed exponent admits only finitely many decompositions into a letter of \(S\) and a finite word in \(E\). There are only finitely many assignments of a word to the finitely many coordinates. This proves both requirements for strong summability.

Ordinary Taylor identities, with legitimate regrouping by the same support bound, prove addition and multiplication. At an ordinary point \(c\), evaluation is \(\sum f_\gamma(c)t^\gamma\). If evaluation vanishes at every ordinary point, uniqueness of Hahn coefficients gives \(f_\gamma(c)=0\) for every \(\gamma,c\). This proves injectivity. Differentiating Taylor coefficients proves the differentiation assertion.

For partial evaluation, derivatives in the parameter variables at \(c\) remain ordinary holomorphic functions on \(V\). The same support calculation applies, independently of the point of \(V\). No shrinking of \(V\) is involved.
\end{proof}

We often suppress the sharp from evaluation notation. Equalities of coherent functions may be checked at all ordinary points. In one variable, equality at ordinary points of a set with an ordinary accumulation point is enough, coefficient by coefficient. This identity statement is distinct from a statement about fine-topological sequences.

\section{A support-increasing operator calculus}\label{sec:operators}
The following lemma packages the common summability argument. The vector spaces may be spaces of holomorphic functions, polynomial remainders, or finite vectors of these.

\begin{lemma}[Strong inverse for a positive operator]\label{lem:operator}
Let \(V\) be a complex vector space and let \(E\subset\Gamma_{>0}\) be well ordered. For each \(e\in E\), let \(T_e:V\to V\) be complex-linear. Extend
\[
 T=\sum_{e\in E}t^eT_e
\]
coefficientwise to Hahn families. For any \(X\in V((t^\Gamma))\), supported in \(S\), the equation
\[
 Y+TY=X
\]
has the unique solution
\begin{equation}\label{eq:operatorinverse}
 Y=\sum_{k\geq0}(-T)^kX,
 \qquad\supp(Y)\subseteq S+E^*.
\end{equation}
If all \(T_e\) preserve holomorphic functions on a fixed domain, every coefficient of \(Y\) is holomorphic on that same domain.
\end{lemma}
\begin{proof}
Expanding \(T^kX\) gives terms indexed by an exponent \(s\in S\) and an ordered word \((e_1,\ldots,e_k)\). Its exponent is \(s+e_1+\cdots+e_k\). The support lemma shows that their union is well ordered and that each exponent receives finitely many contributions, even when all word lengths are allowed. Each coefficient is therefore a finite sum of finite compositions of the \(T_e\) applied to input coefficients. This also proves the fixed-domain assertion.

The products of \(I+T\) with the strongly summable right side telescope coefficientwise, giving \(X\). If \(Z+TZ=0\) with \(Z\ne0\), the least exponent of \(TZ\) is strictly larger than the least exponent of \(Z\). Cancellation at that least exponent is impossible. This proves uniqueness among all Hahn families, not only among those already known to be supported in \(S+E^*\).
\end{proof}

\begin{remark}
There is no assertion here that \(T^kX\) tends to zero in the full fine topology. Even in a fixed ordered value group, ``the valuation increases'' should not be substituted for the finite-word support argument. The latter works without an Archimedean assumption on \(\Gamma\).
\end{remark}

We also use an elementary nonlinear variant. An equation whose coefficient at exponent \(\nu\) is obtained by applying a fixed linear splitting to the input coefficient at \(\nu\), plus finite products of previously determined positive coefficients, can be solved by recursion on \(E^*\). Each recursive product uses exponents strictly smaller than \(\nu\). Well ordering gives the recursion, and finite word multiplicity guarantees that the coefficient formulas are finite. Parameter preparation below is an explicit instance.

\section{Automatic support and a Hahn--Hartogs theorem}\label{sec:hartogs}
\subsection{Pointwise support cannot conceal a decreasing chain}
\begin{lemma}[Automatic common support]\label{lem:autosupport}
Let \(X\) be a connected ordinary complex domain, or a connected complex manifold. Let \((f_\gamma)_{\gamma\in A}\) be a set-indexed family of ordinary holomorphic functions, indexed by a subset \(A\) of an ordered abelian group. Suppose that for every \(x\in X\),
\[
 A_x=\{\gamma\in A:f_\gamma(x)\ne0\}
\]
is well ordered. Then the exact global support
\[
 A'=\{\gamma\in A:f_\gamma\not\equiv0\}
\]
is well ordered.
\end{lemma}
\begin{proof}
If \(A'\) were not well ordered, the ordinary choice principles give a strictly decreasing sequence \(\gamma_1>\gamma_2>\cdots\) in \(A'\). For each \(k\), the zero set of the nonzero holomorphic function \(f_{\gamma_k}\) is closed with empty interior. Its complement \(V_k\) is open and dense. A complex manifold is a Baire space: locally it is a locally compact Hausdorff space. Consequently \(\bigcap_k V_k\) is dense and in particular nonempty. At a point of this intersection all the exponents \(\gamma_k\) occur, contradicting the well ordering of \(A_x\).
\end{proof}

The lemma involves no bound on the number or sizes of the nonzero coefficient functions. Its force comes from the impossibility of making countably many nonzero holomorphic functions vanish collectively at every point.

\subsection{Separate coherence implies joint coherence}
In the next statement a slice is allowed its own workspace and support. The function itself need not initially be presented by any joint series.

\begin{theorem}[Automatic Hahn--Hartogs coherence]\label{thm:hartogs}
Let \(U_1,\ldots,U_d\subseteq\C\) be nonempty connected domains and \(U=\prod_jU_j\). Let
\[
 F:U_1\sh\times\cdots\times U_d\sh\longrightarrow\SC
\]
be a class function. Assume that for every coordinate \(j\), and every fixed choice of the other coordinates in their surcomplex thickenings, the resulting one-variable function is Hahn-coherent on \(U_j\sh\).

Then there exist a set-sized exponent group \(\Gamma\) and a unique joint datum \(H\in\HH(U)\) with \(F=H\sh\). No common-support hypothesis on the slices, no ordinary local boundedness hypothesis on \(F\), and no fine-topological joint continuity hypothesis are required.
\end{theorem}
\begin{proof}
Restrict \(F\) to the ordinary set \(U\). The union
\[
 A=\bigcup_{c\in U}\supp(F(c))
\]
is a set: the ordinary domain is a set, every value has a set-sized support, and class replacement followed by union applies. Let \(\Gamma\) be the additive subgroup generated by \(A\). For \(\gamma\in A\), put
\[
 f_\gamma(c)=[t^\gamma]F(c).
\]
Fix all ordinary coordinates except the \(j\)-th. The coherence of that slice makes \(f_\gamma\) holomorphic in the remaining ordinary coordinate; if \(\gamma\) is absent from the slice support, the coefficient function is zero. Thus each \(f_\gamma\) is separately ordinary holomorphic.

Classical Hartogs separate holomorphy now gives \(f_\gamma\in\OO(U)\). This classical theorem is an input, not a new surcomplex assertion; see, for example, the formulation in \cite{BK}. At any ordinary \(c\), the active coefficient support is exactly \(\supp(F(c))\), which is well ordered. \Cref{lem:autosupport} therefore makes
\(S=\{\gamma:f_\gamma\not\equiv0\}\) well ordered. Define
\[
 H=\sum_{\gamma\in S}f_\gamma t^\gamma\in\HH(U).
\]
By construction \(H\sh=F\) at all ordinary tuples.

We extend this equality one coordinate at a time. Suppose it has been established when the first \(k\) coordinates are arbitrary finite surcomplex points and the remaining ones are ordinary. Fix those first \(k\) points and all ordinary coordinates after \(k+1\). As a function of coordinate \(k+1\), \(F\) is coherent by hypothesis. The function \(H\sh\) on the same slice is coherent by \cref{lem:eval}, after enlarging a workspace if necessary. The two slices agree at every ordinary value of coordinate \(k+1\), by the induction hypothesis. Coefficientwise identity makes them equal on the full thickening of that coordinate. Induction reaches \(k=d\).

Finally, equality of two joint data at all ordinary tuples forces equality of every coefficient function. This proves uniqueness.
\end{proof}

\begin{remark}[The exact role of the slice hypothesis]
The hypothesis concerns fixed \emph{surcomplex} values of all other coordinates, not merely fixed ordinary values. Ordinary-slice data determine the ordinary restriction and its joint Hahn series, but do not by themselves constrain arbitrarily assigned values at tuples with several nonordinary coordinates. The coordinate-by-coordinate argument uses precisely the stronger hypothesis stated in the theorem.
\end{remark}

\begin{corollary}[Coherent Hartogs removal]\label{cor:hole}
Let \(d\geq2\), let \(\Omega\subseteq\C^d\) be a connected domain, and let \(K\subset\Omega\) be compact with \(\Omega\setminus K\) connected. Every Hahn-coherent datum on \((\Omega\setminus K)\sh\) extends uniquely to a Hahn-coherent datum on \(\Omega\sh\), with the same exact support.
\end{corollary}
\begin{proof}
Apply the classical Hartogs extension theorem to each ordinary coefficient function; one formulation is \cite[Theorem 4.3.1]{Lebl}. A nonzero coefficient does not become zero after extension because its restriction is the original coefficient. Thus the same support works. Coefficientwise uniqueness proves the result.
\end{proof}
This corollary is a direct transfer of a classical theorem and is included for completeness. The non-formal issue resolved by \cref{thm:hartogs} is the automatic existence of a common support before a joint datum has been given.

\section{Fixed-domain division with holomorphic parameters}\label{sec:division}
\subsection{An ordinary splitting that does not shrink repeatedly}
Let \(U\subseteq\C^d\) be a connected parameter domain and \(D_r=\{w\in\C:|w|<r\}\). Suppose
\[
 P_0(x,w)=w^m+p_{m-1}(x)w^{m-1}+\cdots+p_0(x)
\]
is ordinary holomorphic in \(x\), and all its roots in \(w\), for every \(x\in U\), lie in \(|w|<r_0<r\).

\begin{lemma}[Uniform ordinary division operators]\label{lem:ordinarydivision}
There are \(\OO(U)\)-linear operators \(D_0,R_0\) on \(\OO(U\times D_r)\), with \(R_0h\in\OO(U)[w]\) of degree less than \(m\), such that
\begin{equation}\label{eq:ordinarysplit}
 h=P_0D_0h+R_0h.
\end{equation}
Both outputs are holomorphic on the same domain \(U\times D_r\). The decomposition is unique, \(D_0(P_0q)=q\), and \(D_0r=0\) for polynomials \(r\) of degree less than \(m\).
\end{lemma}
\begin{proof}
Choose \(r_0<\rho<r\). For \(|w|<\rho\), set
\begin{align}
 D_0h(x,w)
 &=\frac1{2\pi i}\oint_{|\zeta|=\rho}
 \frac{h(x,\zeta)}{P_0(x,\zeta)(\zeta-w)}\,d\zeta,
 \label{eq:D0contour}\\
 R_0h(x,w)
 &=\frac1{2\pi i}\oint_{|\zeta|=\rho}
 \frac{h(x,\zeta)}{P_0(x,\zeta)}
 \frac{P_0(x,\zeta)-P_0(x,w)}{\zeta-w}\,d\zeta.
 \label{eq:R0contour}
\end{align}
The second integrand is polynomial of degree at most \(m-1\) in \(w\), so its integral defines a polynomial for every \(w\). The denominator \(P_0(x,\zeta)\) is nonzero on the integration circle. Differentiation under this ordinary compact contour shows holomorphic dependence on \(x\), locally on \(U\). Combining the two formulas gives the ordinary Cauchy formula and hence \eqref{eq:ordinarysplit} on \(|w|<\rho\).

Outside the roots, \((h-R_0h)/P_0\) extends the first formula to \(D_r\). It is holomorphic near every root because it already agrees there with \eqref{eq:D0contour}; all roots lie inside \(\rho\). Thus the quotient is holomorphic throughout the original disk.

If \(P_0q+r=0\) with \(\deg r<m\), then at each fixed \(x\), the polynomial \(r(x,\cdot)\) vanishes at all roots of \(P_0(x,\cdot)\), with at least their multiplicities. Monicity and degree imply \(r(x,\cdot)=0\), and then \(q=0\). This proves uniqueness, independence of \(\rho\), and the operator identities. The formulas are plainly \(\OO(U)\)-linear.
\end{proof}

For \(P_0=w^m\), these operators are particularly simple:
\begin{equation}\label{eq:monomialoperators}
 R_mh(x,w)=\sum_{j=0}^{m-1}\frac{\partial_w^jh(x,0)}{j!}w^j,
 \qquad
 D_mh=\frac{h-R_mh}{w^m}.
\end{equation}
The apparent singularity in \(D_m\) is removable on the whole disk. The contour construction is the usual analytic division argument, recorded explicitly because its fixed-domain property will be used infinitely often; compare \cite[Theorem 6.2.5]{Lebl}.

\subsection{The matrix theorem}
\begin{theorem}[Support-controlled matrix division]\label{thm:matrixdivision}
Let \(P_0\) satisfy \cref{lem:ordinarydivision}, let \(q\geq1\), and suppose
\[
 \mathcal F=P_0I_q+\mathcal E,
 \qquad \mathcal E\in\Mat_q(\Ip(U\times D_r)).
\]
Let \(E\subset\Gamma_{>0}\) be a well-ordered set containing the supports of all entries of \(\mathcal E\), and put \(M=E^*\). For every vector \(H\in\HH(U\times D_r)^q\), there are unique vectors \(Q\) and \(R\) such that
\begin{equation}\label{eq:matrixdivision}
 H=\mathcal FQ+R,
 \qquad Q\in\HH(U\times D_r)^q,
 \qquad R\in\HH(U)[w]^q,
 \qquad\deg_wR_j<m.
\end{equation}
If \(S\) contains the support of \(H\), then
\[
 \supp(Q)\cup\supp(R)\subseteq S+M.
\]
All coefficients remain holomorphic on \(U\times D_r\). When \(H\) has nonnegative support, so do \(Q\) and \(R\). Consequently
\begin{equation}\label{eq:matrixmodule}
 \Hp(U\times D_r)^q/
 \mathcal F\Hp(U\times D_r)^q
 \cong\Hp(U)^{qm},
\end{equation}
with basis represented by \(w^je_a\), \(0\leq j<m\), \(1\leq a\leq q\). The corresponding assertion holds with \(\HH\) in place of \(\Hp\).
\end{theorem}
\begin{proof}
Apply \(D_0\) to the desired identity, componentwise. Since \(D_0R=0\), it becomes
\[
 Q+D_0(\mathcal E Q)=D_0H.
\]
The operator \(T:Q\mapsto D_0(\mathcal E Q)\) is a positive operator of the form in \cref{lem:operator}. At exponent \(e\), it multiplies by an ordinary holomorphic matrix and then applies the fixed operator \(D_0\). Thus
\begin{equation}\label{eq:matrixQ}
 Q=\sum_{k\geq0}(-D_0\mathcal E)^kD_0H,
 \qquad R=R_0(H-\mathcal E Q),
\end{equation}
where \(D_0\mathcal E\) denotes the operator, not a matrix obtained by differentiating its entries. The support lemma gives the asserted bounds. The ordinary splitting of \(H-\mathcal E Q\) gives
\[
 H-\mathcal E Q=P_0Q+R,
\]
which is \eqref{eq:matrixdivision}.

If two divisions are given, apply \(D_0\) to their difference. The difference of the quotients solves \((I+T)\Delta Q=0\), hence is zero by \cref{lem:operator}. The difference of the remainders is then zero. For nonnegative inputs, \(S+M\subseteq\Gamma_{\geq0}\). Uniqueness says that the bounded-degree vector polynomials are an actual complementary module, not just a spanning set, proving \eqref{eq:matrixmodule}.
\end{proof}

The coefficient matrices in this theorem need not commute with one another. Their finite compositions are ordered words in \eqref{eq:matrixQ}. This observation is important when multiplication by one equation is represented after quotienting by preceding equations.

\subsection{Scalar preparation with parameters}
\begin{theorem}[Fixed-domain Hahn preparation]\label{thm:preparation}
Let \(P_0\) be as above and let
\[
 F=u_0(P_0+E),
 \qquad u_0\in\OO(U\times D_r)^\times,
 \qquad E\in\Ip(U\times D_r).
\]
If \(T\subset\Gamma_{>0}\) contains \(\supp(E)\), then there are unique
\[
 P=P_0+A,\qquad Q=1+B
\]
such that \(F=u_0PQ\), where \(A\in\Ip(U)[w]\) has degree less than \(m\), \(B\in\Ip(U\times D_r)\), and
\[
 \supp(A)\cup\supp(B)\subseteq T^*\setminus\{0\}.
\]
The element \(Q\) is a coherent unit and is nonzero throughout \((U\times D_r)\sh\).
\end{theorem}
\begin{proof}
The equation is
\begin{equation}\label{eq:prepAB}
 E=A+P_0B+AB.
\end{equation}
Recursively through \(T^*\setminus\{0\}\), let
\[
 h_\nu=E_\nu-
 \sum_{\substack{\alpha+\beta=\nu\\\alpha,\beta>0}}A_\alpha B_\beta,
 \qquad A_\nu=R_0h_\nu,
 \qquad B_\nu=D_0h_\nu.
\]
Every sum is finite, and all indices on its right side are smaller than \(\nu\). \Cref{lem:ordinarydivision} verifies \eqref{eq:prepAB} at exponent \(\nu\). It also ensures that every coefficient is holomorphic on the original domain. The positive monoid support makes the completed recursion a Hahn datum.

For uniqueness, compare two normalized pairs and take the least exponent \(\nu\) where either component differs. The product differences in \(AB\) have no contribution there: any contribution contains a lower-index difference. Thus \(\Delta A_\nu+P_0\Delta B_\nu=0\). The uniqueness in ordinary division makes both differences zero, a contradiction. This argument applies to the well-ordered union of the supports of the two candidate pairs.

Finally, \((1+B)^{-1}=\sum_{k\geq0}(-B)^k\) is strongly summable with holomorphic coefficients. Its evaluation is the inverse of \(Q\sh\). Equivalently \(Q\sh\) has standard part one at every finite point.
\end{proof}

\begin{corollary}[Preparation at an ordinary regular point]\label{cor:localprep}
Suppose \(F\in\HH(V)\) has least exponent \(\lambda\), and its leading coefficient \(f_\lambda(x,w)\) is \(w\)-regular of order \(m\) at an ordinary point \((0,0)\). After one ordinary shrink to a product \(U\times D_r\),
\[
 F=t^\lambda u_0PQ
\]
with the factors described in \cref{thm:preparation}. All later Hahn coefficients are constructed on that one product domain.
\end{corollary}
\begin{proof}
Classical Weierstrass preparation gives \(f_\lambda=u_0P_0\) near the point, with \(P_0\) monic of degree \(m\) and \(P_0(0,w)=w^m\). Shrink the ordinary parameter domain and choose concentric fiber disks so that all roots of \(P_0(x,\cdot)\) remain in \(|w|<r_0<r\), while \(u_0\) stays zero-free. The higher Hahn coefficients are already holomorphic on this product, since they were holomorphic on \(V\). Apply \cref{thm:preparation} to \(t^{-\lambda}F/u_0\).
\end{proof}

\subsection{Compatibility with parameter substitution}\label{subsec:basechange}
Let \(\Phi=\phi_0+\eta\) be a finite coherent parameter map from \(V\sh\) to \(U\sh\), with \(\phi_0:V\to U\) ordinary holomorphic and every component of \(\eta\) of positive support. Taylor substitution defines pullback of a datum of support \(S\), with support contained in
\[
 S+\left(\bigcup_j\supp(\eta_j)\right)^*.
\]
The proof is the support calculation of \cref{lem:eval}, with derivatives evaluated at \(\phi_0\).

Normalized preparation and division commute with such pullbacks. Indeed, pulling back their identities gives a normalized factorization or bounded-degree remainder for the pulled-back problem. The relevant uniqueness theorem identifies it with the factorization or remainder constructed directly. For a general \(P_0\), its pullback may itself acquire positive coefficients; the pulled-back leading reference is \(P_0(\phi_0(x),w)\), and those new positive coefficients are incorporated into the perturbation. This is also why specialization at an infinitesimally displaced parameter is legitimate, rather than a purely formal interchange of an infinite tensor product with evaluation.

\section{Coupled systems and finite-free normal forms}\label{sec:finite}
Let
\[
 V=U\times D_{r_1}\times\cdots\times D_{r_n},
 \quad A=\Hp(U),\quad B=\Hp(V),
\]
where \(U\) is a connected ordinary parameter domain; it may also be a point. Let \(\boldsymbol m=(m_1,\ldots,m_n)\) with \(m_i\geq1\), and put
\[
 \boxset=\{\alpha\in\N^n:0\leq\alpha_i<m_i\},
 \qquad N=|\boxset|=\prod_{i=1}^n m_i.
\]
We study arbitrary coupled systems
\begin{equation}\label{eq:system}
 F_i=w_i^{m_i}+E_i,\qquad E_i\in\Ip(V),\quad1\leq i\leq n.
\end{equation}
Let \(E\) be the union of the positive supports of all \(E_i\), and \(M=E^*\). For the unperturbed system take \(E=\varnothing\), so \(M=\{0\}\).

\begin{theorem}[Finite-free complete intersections with controlled support]\label{thm:finitefree}
For \eqref{eq:system}, the following conclusions hold.
\begin{enumerate}
\item The quotient \(\mathcal A=B/(F_1,\ldots,F_n)\) is a free \(A\)-module of rank \(N\), with basis the classes of \(w^\alpha\), \(\alpha\in\boxset\).
\item For every \(H\in\HH(V)\) of support contained in \(S\), there are \(Q_i\in\HH(V)\) and unique \(r_\alpha\in\HH(U)\) such that
\begin{equation}\label{eq:normalform}
 H=\sum_{i=1}^nF_iQ_i+
 \sum_{\alpha\in\boxset}r_\alpha(x)w^\alpha,
\end{equation}
with every output support contained in \(S+M\). All coefficients are holomorphic on the original domains. If \(H\in B\), the outputs may be taken in the corresponding nonnegative-support rings.
\item The sequence \(F_1,\ldots,F_n\) is regular, in any order. In particular, each \(F_i\) is a non-zero-divisor after quotienting by its predecessors, and the final ideal is proper.
\item In the rectangular basis, all structure constants of \(\mathcal A\) have support in \(M\). Their reduction at exponent zero is the multiplication table of
\[
 \OO(U)[w_1,\ldots,w_n]/(w_1^{m_1},\ldots,w_n^{m_n}).
\]
\item The unique normal form commutes with coherent finite parameter substitution. It is independent of the order used to eliminate the equations.
\end{enumerate}
The first, third, and fourth assertions also hold after replacing \(\Hp\) by \(\HH\).
\end{theorem}
\begin{proof}
We prove a stronger inductive statement. After quotienting by the first \(k\) equations, the resulting module is free over
\[
 A_k=\Hp(U\times D_{r_{k+1}}\times\cdots\times D_{r_n}),
\]
with basis \(w_1^{\alpha_1}\cdots w_k^{\alpha_k}\), \(0\leq\alpha_j<m_j\). The normal form has support bound \(S+M\), and every remaining equation acts by a matrix whose reduction is scalar multiplication by its leading coordinate power.

For \(k=1\), apply scalar division, the case \(q=1\) of \cref{thm:matrixdivision}, with \(P_0=w_1^{m_1}\). All variables except \(w_1\) are parameters. This proves the asserted freeness and support estimate.

Suppose the assertion holds after \(k-1\) equations, and put \(q=m_1\cdots m_{k-1}\). Multiplication by \(F_k\) is well defined on the preceding quotient: multiplication by every original function commutes with multiplication by all the defining equations. Represent this action in the established basis. Since the normal form reduces at exponent zero to ordinary division by the preceding coordinate powers, its matrix is
\begin{equation}\label{eq:inducedmatrix}
 \mathcal F_k=w_k^{m_k}I_q+\mathcal E_k,
\end{equation}
where every entry of \(\mathcal E_k\) has positive support contained in \(M\setminus\{0\}\). Its entries are holomorphic on
\(U\times D_{r_k}\times\cdots\times D_{r_n}\).

Apply \cref{thm:matrixdivision} to \eqref{eq:inducedmatrix} in the variable \(w_k\). It splits every vector into an image under multiplication by \(F_k\) and a unique vector polynomial of degree less than \(m_k\). The latter has basis
\[
 w_1^{\alpha_1}\cdots w_{k-1}^{\alpha_{k-1}}w_k^{\alpha_k},
 \qquad 0\leq\alpha_j<m_j,
\]
over \(A_k\). The monoid generated by \(M\setminus\{0\}\) is contained in \(M\), so the new support estimate is still \(S+M\), not a larger uncontrolled support. No domain is shrunk.

At each step the representation of a remaining multiplication operator is obtained by reducing its product with finitely many basis vectors. Hence its entries have the same fixed-domain property and support in \(M\). Their exponent-zero reduction is the expected scalar coordinate power, completing the induction.

For clarity, this quotient construction also gives the actual division identity \eqref{eq:normalform} in the original ring. Each matrix column used above is itself the normal form of multiplication by an original function. Its difference from that original product is a sum of the preceding \(F_i\) times coherent quotients. Substitute these finitely many identities when lifting a vector quotient to the original ring. Their products and sums preserve \(S+M\), since \(M+M=M\). Inducting through the finitely many variables produces the \(Q_i\). The remainders are unique because the rectangular vectors are a free basis. The \(Q_i\) need not be unique, owing to relations among the \(F_i\).

To prove regularity, consider multiplication by \(F_k\) on the preceding free module. A nonzero vector has a least exponent \(\nu\). At that exponent its image under \eqref{eq:inducedmatrix} is \(w_k^{m_k}\) times its nonzero ordinary holomorphic coefficient vector. This vector is nonzero on the connected domain, whereas the positive matrix error cannot contribute at exponent \(\nu\). Thus multiplication by \(F_k\) is injective. The final quotient is nonzero because it is free with the stated basis. The same induction can be performed in any order, proving permutable regularity.

Reducing products of basis monomials proves the structure-constant assertion. Parameter substitution preserves \eqref{eq:normalform} and its bounded degrees, so uniqueness proves its compatibility, as in \cref{subsec:basechange}. Every elimination order gives a normal form in the same rectangular basis of the same quotient; the uniqueness already proved for one order makes all these forms equal. Finally the identical argument works for arbitrary well-ordered input supports, giving the \(\HH\) assertions.
\end{proof}

\subsection{Why this is a finite-mapping theorem}
The algebra \(\mathcal A\) is finite and free over the parameter algebra \(A\). Thus its algebraic projection to the parameter base is finite flat of degree \(N\); equivalently, \(\Spec\mathcal A\to\Spec A\) is a finite flat morphism of rings' spectra. This terminology is algebraic and does not assume that \(A\) is Noetherian. \Cref{thm:fiberlength} below identifies the relevant surcomplex-valued fibers with actual solutions and their local intersection algebras, rather than stopping at a formal module statement.

The hypotheses allow full analytic coupling. For example, the errors can be Hahn sums of functions such as \(\exp(w_1w_2+xw_3)\), with no polynomial bound in the fiber variables. They need not be independent of the coordinate being eliminated. Passing to induced matrices is what makes this possible.

\begin{corollary}[Units and infinitesimal target values]\label{cor:units-targets}
The conclusions of \cref{thm:finitefree} apply after dividing equations whose leading terms are \(u_i(x,w)w_i^{m_i}\), where each \(u_i\) is ordinary holomorphic and zero-free on \(V\). They also apply to
\[
 F_i(x,w)=b_i(x)
\]
for any \(b_i\in\Ip(U)\), by replacing \(E_i\) with \(E_i-b_i\) and using the monoid generated by the enlarged positive support.
\end{corollary}
\begin{proof}
Division by an ordinary zero-free unit changes no exponents and does not change the defining ideal. Subtracting a positive-support target leaves the same coordinate-power leading system. Apply the theorem.
\end{proof}

\subsection{Computational form of the construction}
At an elimination stage, one stores a finite matrix \(\mathcal F_k\) and applies the explicit inverse \eqref{eq:matrixQ}. To compute a specified output exponent, only its finitely many support-word decompositions are needed. This is an exact finite-dependency statement. It is not an assertion that an arbitrary given well-ordered surreal support has a computable enumeration or an efficient representation.

For finitely generated positive supports, one can represent the exponents and enumerate relevant words effectively whenever equality and order in the chosen exponent group are computable. In a usual single-parameter example, the formulas reduce to triangular recursions by powers of that parameter. For a general support, the mathematical recursion is well-founded but may not be an algorithm on finite input. This distinction separates the theorem from a software claim.

\section{Actual fibers and conservation of intersection multiplicity}\label{sec:fibers}
\subsection{Finite Taylor identities at displaced centers}
A possible pitfall is to assume that an arbitrary algebra homomorphism automatically commutes with an infinite Hahn sum. We avoid that assumption by using finite identities in the coherent ring.

\begin{lemma}[Displaced divided differences and finite jets]\label{lem:finitejet}
Let \(D=\prod_{i=1}^nD_{r_i}\), let \(H\) be Hahn-coherent on \(D\), and let \(\xi\in\mm^n\). After choosing a workspace containing the coefficients and \(\xi\), there are coherent functions \(G_j\) on the same \(D\) such that
\begin{equation}\label{eq:divideddifference}
 H(w)-H(\xi)=\sum_{j=1}^n(w_j-\xi_j)G_j(w).
\end{equation}
More generally, for every integer \(\ell\geq1\), there are coherent \(H_\alpha\) on \(D\) such that
\begin{equation}\label{eq:finiteTaylor}
 H(w)=\sum_{|\alpha|<\ell}
 \frac{\partial^\alpha H(\xi)}{\alpha!}(w-\xi)^\alpha
 +\sum_{|\alpha|=\ell}(w-\xi)^\alpha H_\alpha(w).
\end{equation}
Taylor expansion at \(\xi\) is a ring homomorphism to \(K[[u_1,\ldots,u_n]]\), where \(K\) is the chosen coefficient field and \(u=w-\xi\).
\end{lemma}
\begin{proof}
In one variable divide by \(w-\xi\) using \cref{thm:matrixdivision} with \(P_0=w\) and positive error \(-\xi\). The remainder is a coherent function of the other variables. Evaluating at \(w=\xi\), using \cref{lem:eval}, identifies it with \(H(\xi)\). Thus division holds on the original ordinary disk.

Apply this identity successively in all coordinates, telescoping between \(H(w)\) and \(H(\xi)\), to obtain \eqref{eq:divideddifference}. Repeated division by \(w_j-\xi_j\), \(\ell\) times in each variable, gives a rectangular Taylor polynomial and remainders divisible by \((w_j-\xi_j)^\ell\). The coefficients of that polynomial are the indicated derivatives, by repeated differentiation and evaluation. Partition the finitely many polynomial terms of total degree at least \(\ell\) according to a divisor monomial of total degree \(\ell\), and combine them with the existing remainders. This proves \eqref{eq:finiteTaylor}.

Every step is one of the fixed-domain operations already proved. The finite Taylor identities, or the usual Leibniz rule derived from \cref{lem:eval}, show that the Taylor coefficients of a product have the usual finite convolution at each multi-index. Therefore the Taylor map is a ring homomorphism. No convergence in \(u\) is asserted for that formal power series.
\end{proof}

\subsection{Specialization and the finite algebra}
Fix \(a\in U\sh\). Enlarge the exponent group to a divisible set-sized group containing all coefficients and the coordinates of \(a\). Write \(K\) for its algebraically closed Hahn field. By partial evaluation, each \(F_i(a,w)\) is coherent on the same fiber polydisk, with leading term \(w_i^{m_i}\). Applying \cref{thm:finitefree} over this workspace gives
\begin{equation}\label{eq:fiberalgebra}
 \mathcal A_a=
 \mathcal H_K(D)/(F_1(a,w),\ldots,F_n(a,w)),
 \qquad \dim_K\mathcal A_a=N.
\end{equation}
Here \(\mathcal H_K(D)\) means the full coherent ring with exponent group that of \(K\), and constants in \(K\). Its description is still by ordinary complex holomorphic coefficient functions and Hahn support, not by an unspecified analytic structure on \(K\).

Equivalently, evaluate the finitely many multiplication matrices and normal-form coefficients of the parameter algebra. Normal-form uniqueness proves that this gives \eqref{eq:fiberalgebra}. Thus the identification is established directly; it is not based on exchanging evaluation with an infinite coefficient-space tensor product.

\begin{lemma}[All algebraic fiber coordinates are infinitesimal]\label{lem:infcoords}
Every eigenvalue of multiplication by \(w_i\) on \(\mathcal A_a\) is infinitesimal. The characters \(\mathcal A_a\to K\) are in bijection with the actual solutions
\[
 \xi\in D\sh,\qquad F_i(a,\xi)=0\quad(1\leq i\leq n),
\]
and all these solutions lie in \(\mm^n\).
\end{lemma}
\begin{proof}
The multiplication matrix \(W_i\) has nonnegative-valued entries. Its reduction is the nilpotent shift representing \(w_i\) in the truncated polynomial algebra. Hence
\[
 \det(TI-W_i)=T^N+c_{N-1}T^{N-1}+\cdots+c_0,
 \qquad c_j\in\mm.
\]
If a root \(\xi\ne0\) of this polynomial were not infinitesimal, \(\xi^{-1}\) would be finite. Dividing by \(\xi^N\) would express zero as one plus a finite sum of infinitesimals, an impossibility. Zero itself is infinitesimal.

The finite algebra is generated by the coordinate classes, because its rectangular basis consists of coordinate monomials. A character is therefore determined by the tuple \(\xi_i\) to which it sends them. Each \(\xi_i\) is an eigenvalue, for example by the characteristic polynomial identity, and hence infinitesimal. Apply \eqref{eq:divideddifference} to an arbitrary coherent \(H\). Under the character, all factors \(w_j-\xi_j\) vanish. Thus the character sends the class of \(H\) to \(H(\xi)\). In particular \(F_i(a,\xi)=0\).

Conversely, any actual zero gives an evaluation homomorphism annihilating the equations, and hence a character. An actual zero in \(D\sh\) must have standard part satisfying \((\st w_i)^{m_i}=0\), so it lies in \(\mm^n\). This proves both the bijection and the location assertion.
\end{proof}

As a nonzero finite-dimensional commutative algebra over the algebraically closed field \(K\), \(\mathcal A_a\) has finitely many maximal ideals, all with residue field \(K\). Elementary Artinian algebra gives the product decomposition
\begin{equation}\label{eq:artinianproduct}
 \mathcal A_a=\prod_{\xi\in Z_a}\mathcal A_{a,\xi},
\end{equation}
where \(Z_a\) is the finite set of actual roots, and each factor is a local Artinian algebra with nilpotent maximal ideal. In particular, \(|Z_a|\leq N\). We now prove that these factors are exactly the formal analytic intersection algebras.

\subsection{Identification of the local algebras}
\begin{theorem}[Conservation of actual local intersection multiplicity]\label{thm:fiberlength}
For every \(a\in U\sh\) and every actual root \(\xi\in Z_a\), Taylor expansion induces an isomorphism
\begin{equation}\label{eq:localidentification}
 \mathcal A_{a,\xi}\cong
 K[[u_1,\ldots,u_n]]/
 \bigl(F_1(a,\xi+u),\ldots,F_n(a,\xi+u)\bigr).
\end{equation}
The expressions on the right are formal Taylor series of coherent functions. Consequently the local intersection multiplicity
\[
 \mu_\xi(F_a)=
 \dim_K K[[u]]/(F_1(a,\xi+u),\ldots,F_n(a,\xi+u))
\]
is finite and
\begin{equation}\label{eq:conservation}
 \boxed{\displaystyle\sum_{\xi\in Z_a}\mu_\xi(F_a)
 =\prod_{i=1}^n m_i.}
\end{equation}
This count is unchanged by enlarging the divisible Hahn workspace. If the Jacobian determinant \(\det(\partial F_i/\partial w_j)(a,\xi)\) is nonzero at every root, there are exactly \(N\) distinct roots.
\end{theorem}
\begin{proof}
Suppress \(a\) temporarily. The map \(K[w]\to\mathcal A_a\) is surjective, since the rectangular monomials span. Let \(J\) be its kernel, so
\[
 K[w]/J\cong\mathcal A_a.
\]
For \(j\in J\), the division identity \eqref{eq:normalform} has zero remainder, and gives \(j=\sum_iF_iQ_i\) in the coherent ring. Apply the Taylor homomorphism at \(\xi\). If \(\widehat J\) denotes the ideal in \(K[[u]]\) generated by the Taylor polynomials of elements of \(J\), then
\begin{equation}\label{eq:idealfirst}
 \widehat J\subseteq(F_1(\xi+u),\ldots,F_n(\xi+u)).
\end{equation}

We recall explicitly why the standard finite-algebra completion is
\begin{equation}\label{eq:finitecompletion}
 K[[u]]/\widehat J\cong\mathcal A_{a,\xi}.
\end{equation}
Because \(K[w]/J\) is finite-dimensional, its maximal ideals are the finitely many coordinate ideals associated with the tuples in \(Z_a\). In its product decomposition, sufficiently high powers of each local maximal ideal vanish. Localizing \(K[w]/J\) at the ideal of \(\xi\) retains just the factor \(\mathcal A_{a,\xi}\). This factor is already complete for its maximal-ideal filtration, since that filtration terminates. On the polynomial side, localization and then completion gives \(K[[w-\xi]]/\widehat J\). Alternatively, the same isomorphism follows directly by evaluating a formal series in the nilpotent coordinate differences in the local factor: only finitely many terms survive, and a sufficiently high power of the coordinate maximal ideal lies in the localized ideal. These observations prove \eqref{eq:finitecompletion} using only finite-dimensional commutative algebra.

Let \(\ell\) annihilate the maximal ideal of \(\mathcal A_{a,\xi}\). The image of a formal Taylor series in that factor is its polynomial of total degree less than \(\ell\). For \(H=F_i\), the finite identity \eqref{eq:finiteTaylor} shows that this polynomial has the same image as the class of \(F_i\), because every remainder has a factor of total degree \(\ell\) in the nilpotent coordinate differences. That class is zero. Hence \(F_i(\xi+u)\) belongs to \(\widehat J\) for every \(i\). This reverses \eqref{eq:idealfirst} and proves \eqref{eq:localidentification}.

The dimensions of the factors in \eqref{eq:artinianproduct} sum to \(\dim_K\mathcal A_a=N\), giving \eqref{eq:conservation}. If the coefficient workspace is enlarged, the finite matrices extend scalars. The coordinate characteristic polynomials already split over \(K\), so the actual coordinate tuples do not acquire new roots. Each local factor has residue field \(K\) and a nilpotent maximal ideal; after scalar extension it remains local with the same vector-space dimension. Thus the multiplicities are unchanged.

Finally, a nonzero Jacobian determinant is invertible in \(K\). The formal inverse-function theorem, proved by degree-by-degree recursion with invertible linear part, makes the ideal generated by \(F_i(\xi+u)\) equal to the coordinate maximal ideal. The local algebra is then \(K\), of dimension one. If this holds at every root, \eqref{eq:conservation} gives \(|Z_a|=N\).
\end{proof}

\begin{remark}[Formal local versus fine analytic language]
The definition of \(\mu_\xi\) is formal and unambiguous. The coherent Taylor series also represents the function at sufficiently small surcomplex displacements, by the infinitesimal evaluation theory of \cite{Source}. Thus \eqref{eq:localidentification} relates actual surcomplex zeros to their analytic jets. It does not assume a general theorem identifying arbitrary fine-analytic sheaves with Noetherian convergent-series rings.
\end{remark}

\begin{corollary}[Infinitesimal conservation law]
Every coupled positive-support perturbation of the coordinate-power system has the same total local intersection multiplicity \(N\) in the whole infinitesimal fiber monad, at every admissible parameter. The roots may split into smaller subclusters, lie at different infinitesimal scales, or collide. None of these changes alters the total in \eqref{eq:conservation}.
\end{corollary}

\section{A multidimensional coherent residue functional}\label{sec:residue}
\subsection{Definition without fine-continuous contours}
Retain \eqref{eq:system}. Choose ordinary radii \(0<\rho_i<r_i\) and let \(T_\rho\) denote the product of the positively oriented ordinary circles \(|w_i|=\rho_i\). On an ordinary product annulus around this torus, expand
\begin{equation}\label{eq:inverseFi}
 \frac1{F_i}
 =\sum_{k_i\geq0}(-1)^{k_i}
 \frac{E_i^{k_i}}{w_i^{m_i(k_i+1)}}.
\end{equation}
The coefficient functions are holomorphic on that annulus, and the Hahn supports are positive-monoid controlled. For \(H\in\HH(V)\), define
\begin{equation}\label{eq:lambda}
 \Lambda_F(H)
 =\frac1{(2\pi i)^n}\hoint_{T_\rho}
 \frac{H(x,w)\,dw_1\cdots dw_n}{F_1(x,w)\cdots F_n(x,w)}.
\end{equation}
The superscript \(\mathrm H\) instructs us to integrate each ordinary complex coefficient and then assemble the Hahn series. It does not denote integration along a fine-continuous path in \(\SC\).

\begin{theorem}[Existence, coefficient formula, and ideal annihilation]\label{thm:residuefunctional}
The expression \eqref{eq:lambda} is well defined, independent of the radii and of the order of iterated integration. It is \(\HH(U)\)-linear in \(H\), with
\[
 \supp\Lambda_F(H)\subseteq S+M
 \quad\text{if}\quad\supp(H)\subseteq S.
\]
It is given explicitly by
\begin{equation}\label{eq:coefficientresidue}
 \Lambda_F(H)=
 \sum_{k\in\N^n}(-1)^{|k|}
 \left[w_1^{m_1(k_1+1)-1}\cdots
       w_n^{m_n(k_n+1)-1}\right]
 \left(H\prod_{i=1}^nE_i^{k_i}\right).
\end{equation}
Here coefficient extraction is the ordinary Taylor coefficient in \(w\), extended coefficientwise in the Hahn exponents. The sum is strongly summable. Moreover
\[
 \Lambda_F(F_jH)=0\qquad(1\leq j\leq n).
\]
Thus \(\Lambda_F\) descends to the finite quotient algebra, and it maps its integral version to \(\Hp(U)\). All constructions commute with coherent finite parameter substitution.
\end{theorem}
\begin{proof}
Expand the product of \eqref{eq:inverseFi}. If the numerator has support \(S\), the complete expansion has support in \(S+M\), and a fixed exponent has only finitely many contributions, by the finite-word form of the support lemma. Thus at each exponent the proposed contour integrand is a finite sum of ordinary holomorphic functions on the product annulus.

For a fixed tuple \(k\), the ordinary repeated Cauchy coefficient formula gives precisely the summand in \eqref{eq:coefficientresidue}. This proves the formula and independence of all radii. Ordinary integration of each coefficient preserves holomorphic parameter dependence on \(U\); it can be checked locally by differentiation under the compact torus. Ordinary Fubini for continuous integrands on the product of circles proves independence of order. All these operations are performed on finite sums at each Hahn exponent, so no exchange of non-absolutely-convergent ordinary infinite sums is being used.

To prove annihilation, cancel \(F_j\) in \(F_jH/(F_1\cdots F_n)\), using the strongly summable geometric inverse identity on the annulus. In the remaining expansion, denominator monomials involve only \(w_i\) for \(i\ne j\). Each coefficient of the numerator is holomorphic in \(w_j\) on the entire disk. Its integral in \(w_j\) is zero. Hence every Hahn coefficient of \(\Lambda_F(F_jH)\) vanishes.

Linearity follows coefficientwise. Nonnegative input support remains nonnegative after adding \(M\). Finally, ordinary contour integration commutes with parameter derivatives on a fixed torus. Taylor substitution in a positive parameter displacement is strongly summable with the bounds already proved. Performing it before or after the contour functional therefore gives the same result.
\end{proof}

\subsection{Perfect pairing through collisions}
\begin{theorem}[Relative residue duality]\label{thm:duality}
Let \(\mathcal A=B/(F_1,\ldots,F_n)\) and \(A=\Hp(U)\) as in \cref{thm:finitefree}. The map
\begin{equation}\label{eq:dualmap}
 \mathcal A\longrightarrow\Hom_A(\mathcal A,A),
 \qquad g\longmapsto\bigl(h\longmapsto\Lambda_F(gh)\bigr)
\end{equation}
is an isomorphism. In the rectangular basis \(b_\alpha=w^\alpha\), the Gram matrix
\[
 G_{\alpha\beta}=\Lambda_F(w^{\alpha+\beta})
\]
has support in \(M\) and satisfies
\begin{equation}\label{eq:Gramleading}
 (G_0)_{\alpha\beta}=
 \begin{cases}
 1,&\alpha_i+\beta_i=m_i-1\text{ for every }i,\\
 0,&\text{otherwise}.
 \end{cases}
\end{equation}
Its inverse has support in \(M\). The same pairing is perfect over the full parameter ring \(\HH(U)\), and after specialization at every surcomplex parameter.
\end{theorem}
\begin{proof}
At exponent zero, every error in \eqref{eq:coefficientresidue} disappears. The functional is the coefficient of \(w_1^{m_1-1}\cdots w_n^{m_n-1}\). This gives \eqref{eq:Gramleading}. The map
\(\alpha\mapsto(m_1-1-\alpha_1,\ldots,m_n-1-\alpha_n)\)
is an involution of the rectangular index set. Thus \(G_0\) is a permutation matrix and \(G_0^2=I\).

Write \(G=G_0+C\), where \(C\) has positive support in \(M\). Then
\begin{equation}\label{eq:Graminverse}
 G^{-1}=\sum_{k\geq0}(-G_0C)^kG_0.
\end{equation}
This matrix series is strongly summable by \cref{lem:operator}, and its support remains in \(M\). Hence \(G\) is invertible over \(A\). Since \(\mathcal A\) is finite free, invertibility of this Gram matrix is exactly the isomorphism \eqref{eq:dualmap}.

The same matrix is invertible after adjoining arbitrary monomials or after any allowed parameter evaluation. Positivity of its nonconstant coefficients is preserved by finite coherent substitution. This proves the remaining assertions.
\end{proof}

This is a finite-free Frobenius algebra structure: the algebra is identified with its parameter-linear dual by one explicitly defined functional. We use this description rather than silently assuming a Noetherian version of relative Gorenstein duality. The proof shows directly what remains valid over the present coefficient ring.

Most importantly, \(\det G\) has nonzero \emph{constant} leading coefficient. No inverse discriminant is needed. A collision of roots can destroy a decomposition into distinct point evaluations, but it cannot make this residue pairing degenerate.

\subsection{Local summands and a canonical trace element}
Specialize to a field \(K\) and use \eqref{eq:artinianproduct}. Let \(e_\xi\) be the identity idempotent of the factor at \(\xi\). Define
\[
 \Lambda_{F,\xi}(H)=\Lambda_F(e_\xi H).
\]
Then
\begin{equation}\label{eq:localresiduesum}
 \Lambda_F(H)=\sum_{\xi\in Z_a}\Lambda_{F,\xi}(H).
\end{equation}
The summands are finite in number. Factors at distinct roots are orthogonal for the pairing because their product is zero. The restricted pairing on each local factor is perfect: an element annihilating that factor also pairs to zero with every other factor and hence is zero. This gives localized residue functionals on the actual local intersection algebras of \cref{thm:fiberlength}.

For completeness, the following purely finite-algebra identity extracts traces without requiring a formula for individual residue weights.

\begin{proposition}[Trace element from the residue pairing]\label{prop:Euler}
Let \(b_\alpha=w^\alpha\), and let \(b^\alpha\) be the residue-dual basis, so that \(\Lambda_F(b_\beta b^\alpha)=\delta_{\alpha\beta}\). Define
\begin{equation}\label{eq:Eulerelement}
 \mathfrak e_F=\sum_{\alpha\in\boxset}b_\alpha b^\alpha\in\mathcal A.
\end{equation}
For every \(H\in\mathcal A\),
\begin{equation}\label{eq:traceEuler}
 \Tr(m_H)=\Lambda_F(H\mathfrak e_F),
 \qquad
 (\mathfrak e_F)_0=Nw_1^{m_1-1}\cdots w_n^{m_n-1}.
\end{equation}
Here \(m_H\) is multiplication by \(H\). The element \(\mathfrak e_F\) is independent of the chosen basis; its normal-form coefficients have support in \(M\). After specialization,
\begin{equation}\label{eq:tracelength}
 \Tr(m_H)=\sum_{\xi\in Z_a}\mu_\xi(F_a)H(\xi).
\end{equation}
\end{proposition}
\begin{proof}
The coefficient of \(b_\alpha\) in the expansion of \(Hb_\alpha\) is \(\Lambda_F(Hb_\alpha b^\alpha)\). Summing these diagonal coefficients gives the first identity in \eqref{eq:traceEuler}. The functional \(H\mapsto\Tr(m_H)\) is basis independent, and the perfect pairing identifies it with a unique element; hence \(\mathfrak e_F\) is basis independent.

The inverse Gram matrix gives the dual basis, so \cref{thm:duality} and the multiplication-table support estimate place \(\mathfrak e_F\) in support \(M\). At exponent zero, the dual of \(w^\alpha\) is \(w^{\boldsymbol m-\boldsymbol1-\alpha}\). Each product is the top rectangular monomial, and there are \(N\) terms. This proves the leading formula.

On the local factor at \(\xi\), an element \(H\) is \(H(\xi)\) plus an element of the nilpotent maximal ideal. Multiplication by the latter is nilpotent and has trace zero. Multiplication by the former scalar has trace \(\dim_K\mathcal A_{a,\xi}\,H(\xi)\). Sum over the local factors and use \cref{thm:fiberlength}.
\end{proof}

\begin{remark}[Normalization and comparison with other residue theories]\label{rem:residuecomparison}
The functional \(\Lambda_F\) is specified by \eqref{eq:lambda}, including the order of the defining equations and the orientation convention. Its local summands are defined by algebraic idempotents and are proved nondegenerate. We do not assume without proof a universal identification with every Grothendieck-residue construction over non-Noetherian Hahn analytic rings. Nor do we need a general identity \(\mathfrak e_F=\det(\partial F_i/\partial w_j)\) to prove duality or trace conservation. In one variable we prove the familiar derivative formula below, and in the coupled example we verify the Jacobian weights directly. Extending that identification in complete generality is a compatibility question beyond the theorems asserted here, not a claim that the classical identity is an open problem.
\end{remark}

\section{One-variable collision invariants and lifted monodromy}\label{sec:onevar}
This section gives two complementary consequences of parameter division. The one-variable residue identity itself is already present without this parameter organization in \cite{Source}; the point here is its common-domain, collision-stable form.

\subsection{Remainders, residues, and traces}
Let \(P\in\Hp(U)[w]\) be monic of degree \(m\), with leading ordinary polynomial \(P_0\) satisfying the root-disk condition of \cref{lem:ordinarydivision}, and \(P-P_0\) of positive support. Divide \(H\in\HH(U\times D_r)\) by \(P\), with unique remainder \(R_H\) of degree less than \(m\). Define
\[
 \lambda_P(H)=\frac1{2\pi i}\hoint_{|w|=\rho}\frac{H}{P}\,dw,
 \qquad r_0<\rho<r.
\]
The inverse of \(P\) on this ordinary circle is obtained by geometric expansion about \(P_0\).

\begin{proposition}[One-variable trace and residue formulas]\label{prop:univtrace}
For the rank-\(m\) algebra \(\mathcal A_P=\Hp(U\times D_r)/(P)\),
\begin{align}
 \lambda_P(H)&=[w^{m-1}]R_H,\label{eq:univresidue}\\
 \Tr(m_H)&=\lambda_P(P_wH)
 =\frac1{2\pi i}\hoint_{|w|=\rho}\frac{HP_w}{P}\,dw.
 \label{eq:univtrace}
\end{align}
The pairing \((G,H)\mapsto\lambda_P(GH)\) is perfect over the parameter ring. All these quantities are coherent across the ordinary discriminant locus, and specialize to the corresponding formulas for the actual surcomplex roots with their multiplicities.
\end{proposition}
\begin{proof}
The quotient term \(H/P-R_H/P\) is coherent holomorphic on the disk, so its contour integral vanishes. At each Hahn exponent, the expansion of \(R_H/P\) is an ordinary rational function with possible poles only at roots of \(P_0\), all inside the circle. Its contour integral equals its coefficient of \(w^{-1}\) at infinity. Since \(P\) is monic and \(\deg R_H<m\), that coefficient is exactly \([w^{m-1}]R_H\). The claim is coefficientwise and does not require extending \(H\) outside its original disk.

For monomials \(w^i,w^j\), the Gram entry is zero when \(i+j<m-1\), and is one when \(i+j=m-1\). Reverse the order of one basis index; the matrix is triangular with diagonal ones. Its determinant is \((-1)^{m(m-1)/2}\), independently of the coefficients of \(P\). Thus the pairing is perfect even when \(P\) has repeated roots.

To establish the trace formula, first specialize an ordinary parameter \(x\) and choose an algebraically closed Hahn workspace containing the coefficients. Factor \(P=\prod_\xi(w-\xi)^{d_\xi}\). All roots are finite with standard parts among the roots of \(P_0\). Polynomial division and ordinary local preparation identify the quotient with the product of the local algebras \(K[u]/(u^{d_\xi})\); the same conclusion follows by factoring the monic polynomial and using finite Taylor remainders. Multiplication by \(H\) has trace \(\sum_\xi d_\xi H(\xi)\).

On the other hand \(P_w/P=\sum_\xi d_\xi/(w-\xi)\). For a finite root \(\xi=c+\varepsilon\), expand \((w-\xi)^{-1}=\sum_{k\geq0}\varepsilon^k(w-c)^{-k-1}\) on the ordinary circle. Ordinary Cauchy derivative formulas, applied to each Hahn coefficient, give the value \(H(\xi)\). The support is the numerator support plus \(\supp(\varepsilon)^*\), so this computation is strongly summable. Thus the contour integral is \(\sum_\xi d_\xi H(\xi)\). Thus \eqref{eq:univtrace} holds at every ordinary parameter. The quantities on both sides are coherent functions of the parameter, so coefficient uniqueness proves the identity on \(U\), and allowed specialization gives it at every finite surcomplex parameter.

Division, integration, and finite matrix traces all have holomorphic parameter coefficients on \(U\). None uses an inverse discriminant. This proves the collision statement.
\end{proof}

\subsection{Monodromy above the ordinary simple-root locus}
Suppose \(P_0(x,\cdot)\) has only simple roots for \(x\in U\). Its ordinary root space
\[
 X_0=\{(x,c)\in U\times D_r:P_0(x,c)=0\}
\]
is an unramified finite analytic covering of \(U\). A coherent function on this ordinary manifold means a family of ordinary holomorphic coefficient functions in its charts, with a common well-ordered support; equality on chart overlaps is coefficientwise.

\begin{theorem}[Support-preserving lift of the ordinary root cover]\label{thm:monodromy}
Let \(P=P_0+E\) as above and suppose \(P_0\) is simple-rooted over \(U\). There is a unique coherent function on \(X_0\) of the form
\[
 r(x,c)=c+\eta(x,c),\qquad\supp(\eta)\subseteq\supp(E)^*\setminus\{0\},
\]
satisfying \(P(x,r(x,c))=0\). On every root chart, its coefficients are holomorphic on that chart's entire ordinary domain. Analytic continuation of the lifted roots along an ordinary loop permutes their labels in exactly the same way as continuation of the ordinary roots of \(P_0\).
\end{theorem}
\begin{proof}
On \(X_0\), the function \(P_{0,w}(x,c)\) is ordinary holomorphic and nowhere zero. Substituting \(c+\eta\) into \(P\), the coefficient of \(\eta_\nu\) at a positive exponent \(\nu\) is \(P_{0,w}(x,c)\eta_\nu\). All other terms at that exponent involve the input and previously determined coefficients: the nonlinear terms contain at least two positive corrections, and a correction in \(E\) is multiplied by a positive input exponent. Recursively divide their negative sum by \(P_{0,w}(x,c)\). The support-word argument constructs \(\eta\) on \(\supp(E)^*\), with finite contributions at each exponent.

Every operation uses ordinary holomorphic functions on \(X_0\) and division by the globally zero-free derivative. The coefficients therefore glue on the root space itself; no choice of globally labeled roots over \(U\) is needed. Uniqueness follows from the least positive exponent at which two candidate corrections differ. Continuing a chart expression along a loop gives the expression on the continued ordinary sheet, and uniqueness excludes any additional permutation. The resulting monodromy is exactly that of the ordinary cover.
\end{proof}

The loops here live in the ordinary parameter base, and continuation concerns ordinary holomorphic coefficients. This theorem does not postulate nonconstant fine-continuous real-parameter loops in the surcomplex class. New collisions whose parameter standard part lies on the ordinary discriminant are also not excluded; the theorem deliberately works over its complement.

\section{Worked examples}\label{sec:examples}
\subsection{A genuinely coupled four-point splitting}\label{subsec:fourpoints}
Let \(t=\omega^{-1}\) and consider
\begin{equation}\label{eq:fourpointssystem}
 F_1=x^2-ty,\qquad F_2=y^2-tx.
\end{equation}
Here \(x,y\) are fiber variables; there are no external parameters. \Cref{thm:finitefree} gives the basis \((1,x,y,xy)\) and dimension four. The quotient relations are
\[
 x^2=ty,\qquad y^2=tx,\qquad
 x^2y=t^2x,\qquad xy^2=t^2y,\qquad x^2y^2=t^2xy.
\]
In the indicated basis, multiplication by \(x\) and \(y\) is represented by
\begin{equation}\label{eq:fourpointsmatrices}
 W_x=\begin{pmatrix}
 0&0&0&0\\1&0&0&t^2\\0&t&0&0\\0&0&1&0
 \end{pmatrix},\qquad
 W_y=\begin{pmatrix}
 0&0&0&0\\0&0&t&0\\1&0&0&t^2\\0&1&0&0
 \end{pmatrix}.
\end{equation}
They commute and satisfy \(W_x^2=tW_y\), \(W_y^2=tW_x\). Both characteristic polynomials are \(T^4-t^3T\).

The actual roots are
\begin{equation}\label{eq:fourpointsroots}
 (0,0),\qquad(t\zeta,t\zeta^2)\quad(\zeta^3=1).
\end{equation}
Indeed, if \(x=0\), the second equation gives \(y=0\). Otherwise \(y=x^2/t\), and substitution gives \(x^3=t^3\). The Jacobian determinant is
\[
 J=\det\begin{pmatrix}2x&-t\\-t&2y\end{pmatrix}=4xy-t^2.
\]
It equals \(-t^2\) at the origin and \(3t^2\) at each of the other roots. Thus all four roots are simple, in agreement with the total multiplicity theorem. The coupling is essential: neither coordinate equation alone determines a finite set in two variables.

The coefficient formula \eqref{eq:coefficientresidue} yields
\[
 \Lambda_F(1)=\Lambda_F(x)=\Lambda_F(y)=0,
 \qquad \Lambda_F(xy)=1.
\]
For example, a contribution to \(\Lambda_F(xy)\) would require
\(k_2=2k_1\) and \(k_1=2k_2\), so only \(k_1=k_2=0\) contributes. For the other three basis elements the corresponding exponent equations have no nonnegative solution. Therefore
\begin{equation}\label{eq:fourpointsGram}
 G=\begin{pmatrix}
 0&0&0&1\\
 0&0&1&0\\
 0&1&0&0\\
 1&0&0&t^2
 \end{pmatrix},\qquad\det G=1.
\end{equation}
The dual basis is \((xy-t^2,y,x,1)\), and consequently
\[
 \mathfrak e_F=(xy-t^2)+xy+xy+xy=4xy-t^2=J.
\]
Here the trace element agrees with the Jacobian by a direct calculation, not an unproved general identification.

For every coherent observable \(H\), we have the exact root-weight formula
\begin{equation}\label{eq:fourpointsresidues}
 \boxed{\displaystyle
 \Lambda_F(H)=
 -\frac{H(0,0)}{t^2}
 +\frac1{3t^2}\sum_{\zeta^3=1}H(t\zeta,t\zeta^2).}
\end{equation}
To prove it, the right side annihilates the defining equations because it is a linear combination of evaluations at roots. Check equality with \(\Lambda_F\) on the four basis monomials, using \(\sum\zeta=\sum\zeta^2=0\); normal-form uniqueness then proves equality for every \(H\). Thus in this example the local weights are exactly \(1/J(\xi)\).

The example also exhibits a specifically non-Archimedean cancellation. Each nonzero weight in \eqref{eq:fourpointsresidues} is infinite, but for \(H=1\) their sum is exactly zero. The globally defined residue pairing has nonnegative support and determinant one. Splitting it into individual root weights can introduce negative exponents that cancel only after the finite sum is taken.

\subsection{The same multiplicity for a nonpolynomial coupled system}\label{subsec:sin}
On ordinary disks sufficiently small that \(\sin z/z\) is zero-free, consider
\begin{equation}\label{eq:sinsystem}
 \sin^2x-ty=0,\qquad \sin^2y-tx=0.
\end{equation}
The squares of \(\sin x/x\) and \(\sin y/y\) are ordinary units. Dividing the equations by them normalizes the leading terms to \(x^2,y^2\). \Cref{cor:units-targets,thm:fiberlength} therefore gives total intersection multiplicity four in \(\mm^2\), although both equations have infinitely many ordinary Taylor terms and are coupled.

The origin is an exact solution, with Jacobian determinant \(-t^2\). For the other solutions set \(x=tX\), \(y=tY\), and \(s=t^2\). The rescaled equations have ordinary formal expansions
\[
 X^2-Y-\frac{sX^4}{3}+O(s^2)=0,
 \qquad
 Y^2-X-\frac{sY^4}{3}+O(s^2)=0.
\]
At \(s=0\), the three nonzero solutions are \((X,Y)=(\zeta,\zeta^2)\), \(\zeta^3=1\), with Jacobian determinant three. The ordinary formal implicit-function recursion gives a unique power series in \(s\) at each such point. Its evaluation at the infinitesimal \(s\) is strongly summable.

Writing \(X=\zeta+As+O(s^2)\), \(Y=\zeta^2+Bs+O(s^2)\), the first correction equations are
\[
 2\zeta A-B=\frac\zeta3,\qquad
 2\zeta^2B-A=\frac{\zeta^2}3.
\]
Thus the three nonzero roots are
\begin{align}
 x_\zeta&=t\zeta+\frac{2+\zeta^2}{9}t^3+O(t^5),\\
 y_\zeta&=t\zeta^2+\frac{2+\zeta}{9}t^3+O(t^5).
 \label{eq:sineroots}
\end{align}
Here \(O(t^5)\) denotes a formal power series in \(t\) beginning at degree five, evaluated strongly. The Jacobian determinant at each nonzero root is \(3t^2+O(t^4)\), so all three are simple. Along with the origin they already account for the full multiplicity four, proving there are no hidden additional roots at smaller or larger infinitesimal subscales in this monad.

\subsection{A support with infinitely many exponents below one}\label{subsec:nondiscrete}
Let
\[
 \gamma_k=1-\frac1{k+2},\qquad k=0,1,2,\ldots,
\]
and consider, on any fixed ordinary product of disks,
\begin{equation}\label{eq:nondiscreteexample}
 F_1(x,y)=x^2+\sum_{k\geq0}t^{\gamma_k}e^{k!x+y},
 \qquad
 F_2(x,y)=y^3+t e^{x-y}.
\end{equation}
The support \(\{\gamma_k:k\geq0\}\) is well ordered of order type \(\omega\), even though it has supremum one. The common error support includes one and has order type \(\omega+1\). Its coefficients are entire ordinary functions, with no uniform growth restriction as \(k\) increases.

The theorem applies with \(m_1=2,m_2=3\). The quotient has the six-element basis
\[
 1,\ y,\ y^2,\ x,\ xy,\ xy^2,
\]
and every actual fiber has total intersection multiplicity six. All normal-form and residue coefficients are controlled by
\[
 M=\left(\{\gamma_k:k\geq0\}\cup\{1\}\right)^*.
\]
There are infinitely many input terms below the exponent one. At \(x=y=0\), replacing \(t\) by any fixed ordinary real \(q\in(0,1)\) would make the summands \(q^{\gamma_k}\) tend to \(q\), not to zero. Thus ordinary numerical convergence of this perturbation series is unavailable even at that point. It also cannot be an ordinary integer-power series in one positive monomial: there are infinitely many distinct exponents in a bounded interval. The strong-support proof, however, needs neither of these properties.

This example is useful for locating the real generality of the finite-mapping theorem. It is not just a restatement for a discrete one-parameter power-series family in which analytic convergence happens to have been omitted.

\subsection{Residue duality versus trace-pairing degeneration}\label{subsec:cubic}
Let \(a,b\) be positive-support coherent parameter functions and put
\[
 P(w)=w^3-aw-b.
\]
The basis is \((1,w,w^2)\), and \(w^3=aw+b\). By \cref{prop:univtrace}, the residue Gram matrix is
\begin{equation}\label{eq:cubicGram}
 G_{\mathrm{res}}=
 \begin{pmatrix}0&0&1\\0&1&0\\1&0&a\end{pmatrix},
 \qquad \det G_{\mathrm{res}}=-1.
\end{equation}
The residue-dual basis is \((w^2-a,w,1)\), giving
\[
 \mathfrak e_P=3w^2-a=P_w.
\]
By contrast, the matrix of the \emph{trace pairing} \((H,G)\mapsto\Tr(m_{HG})\) is
\begin{equation}\label{eq:cubictraceGram}
 G_{\mathrm{tr}}=
 \begin{pmatrix}
 3&0&2a\\0&2a&3b\\2a&3b&2a^2
 \end{pmatrix},
 \qquad \det G_{\mathrm{tr}}=4a^3-27b^2.
\end{equation}
At a repeated root the trace pairing can degenerate; the residue pairing does not. This distinction is essential in using the word ``duality.''

For an explicit collision take \(a=3t^2\), \(b=2t^3\). Then
\[
 P=(w-2t)(w+t)^2,
\]
and the actual local multiplicities are one at \(2t\) and two at \(-t\). A partial-fraction calculation gives
\begin{equation}\label{eq:doublepolelambda}
 \lambda_P(H)=
 \frac{H(2t)-H(-t)}{9t^2}-\frac{H'(-t)}{3t},
\end{equation}
while the trace is
\[
 \Tr(m_H)=H(2t)+2H(-t).
\]
Formula \eqref{eq:doublepolelambda} follows from
\[
 \frac1P=\frac1{9t^2(w-2t)}-
 \frac1{9t^2(w+t)}-\frac1{3t(w+t)^2}
\]
and the one-variable Cauchy derivative formula. It shows concretely how the nondegenerate functional remembers a first derivative at the double point, information invisible to the trace pairing on nilpotent directions.

\subsection{Ordinary monodromy and new infinitesimal collisions}
Consider
\[
 P(z,w)=w^3-z-tw.
\]
On an ordinary annulus \(0<r<|z|<R\), choose a sufficiently large ordinary fiber disk. The leading polynomial \(w^3-z\) has three simple roots, with the ordinary cyclic monodromy around zero. \Cref{thm:monodromy} lifts these sheets coherently and preserves that three-cycle. On a local choice \(c^3=z\), the first correction is
\[
 w=c+\frac{t}{3c}+\text{terms of higher positive support}.
\]
The exact polynomial discriminant is \(4t^3-27z^2\), so collisions occur at
\[
 z=\pm\frac{2}{3\sqrt3}t^{3/2}.
\]
Both have standard part zero and therefore lie outside the thickening of the ordinary annulus. The monodromy theorem neither misses nor forbids them: its domain intentionally excludes the whole standard-zero parameter monad. Meanwhile the finite algebra and residue pairing remain meaningful across such collisions on a larger parameter chart to which preparation applies.

\section{Relation to published work and novelty audit}\label{sec:literature}
\subsection{What is inherited}
Alling's book \cite{Alling} and the work of van den Dries and Ehrlich \cite{vdDE} are part of the established surreal analytic background. Hahn-field algebra and algebraic closure are classical; the specific algebraic-closure input used here is \cite[Corollary 4]{Poonen}. The ordinary Hartogs and Weierstrass theorems are classical complex analysis, not new assertions of this article. The supplied manuscript \cite{Source} already develops strong evaluation, one-variable coherent preparation, exact root lifting, and coherent residues at moving poles. Our use of those concepts is an extension of that manuscript, not a claim to have originated them.

Berarducci and Mantova's work on transseries as germs of surreal functions \cite{BM} develops composition and Taylor behavior for a distinguished surreal-function class. The present coefficient derivative differentiates ordinary complex variables and leaves the monomials \(t^\gamma\) fixed. It should not be confused with a derivation of the surreal field itself.

\subsection{Why existing non-Archimedean division must be acknowledged}
Cluckers and Lipshitz give a broad theory of fields with analytic structure \cite{CL}, and their later paper treats strictly convergent analytic structures and establishes, in considerable generality, implications from Weierstrass preparation to division \cite{CLstrict}. Consequently ``a non-Archimedean Weierstrass theorem'' is not, by itself, a justified novelty claim.

Our precise assertions concern ordinary holomorphic coefficient functions on a fixed ordinary domain, one global well-ordered Hahn support, no ordinary norm convergence of the Hahn series, and the explicit bounds \(S+E^*\). The matrix induction retains those bounds for coupled systems and then identifies actual surcomplex formal local algebras. The proofs here do not require that this coefficient category be installed as one of the existing analytic structures. Conversely, we have not proved that every theorem of those frameworks transfers to it, or that no existing general theorem subsumes the present results. A complete comparison would require matching the exact rings, admissibility axioms, and function interpretations.

\subsection{A terminology warning about Hahn-holomorphic functions}
M\"uller and Strohmaier's \emph{The theory of Hahn-meromorphic functions, a holomorphic Fredholm theorem, and its applications} \cite{MS} studies functions near zero on the logarithmic cover of the ordinary complex plane, using convergent generalized-series expansions. Their setup includes generalized powers and logarithmic terms and imposes analytic convergence and admissibility conditions.

The Hahn-coherent objects used here are different: the variable \(w\) is an ordinary holomorphic coefficient variable with an infinitesimal thickening, while \(t^\gamma\) records surreal Hahn scales. Strong summability permits supports and coefficient growth that need not give a convergent ordinary series after replacing \(t\) by a positive real number, as \cref{subsec:nondiscrete} demonstrates. The shared word ``Hahn'' does not identify the two function categories. We use ``Hahn-coherent'' consistently to preserve the terminology of the supplied manuscript and avoid claiming its category is identical to that published theory.

\subsection{Contribution-by-contribution assessment}
\begin{longtable}{@{}L{0.29\textwidth}L{0.62\textwidth}@{}}
\toprule
Result & Assessment and source of the proof\\
\midrule
\endhead
Automatic support and separate coherence & The Baire support lemma removes a global-support assumption. It is combined with classical Hartogs to prove the exact slice theorem in this category. This formulation was not located in the targeted search.\\[0.5em]
Fixed-domain matrix division & A strong Neumann inverse built on ordinary division. The matrix and fixed-domain support bookkeeping is explicit; general non-Archimedean division has extensive predecessors.\\[0.5em]
Coupled finite-free algebra & Iterated induced matrix division, not independent scalar root counting. The result includes arbitrary holomorphic cross-couplings and nondiscrete positive Hahn supports.\\[0.5em]
Actual multiplicity conservation & The finite algebra is identified with the formal local algebras of its actual infinitesimal roots. Finite displaced Taylor identities avoid unjustified interchange by abstract homomorphisms.\\[0.5em]
Relative residue duality & A coefficientwise multidimensional functional has an invertible Gram matrix with explicitly known reduction. Classical residue duality is the analogue; the arbitrary-support, fixed-domain construction is proved directly.\\[0.5em]
One-variable traces and monodromy & Parameter-organized consequences of division, classical analytic continuation, and the established one-variable contour framework. These are not asserted to be wholly new general principles.\\
\bottomrule
\end{longtable}

The literature inspection used primary publisher pages, author-hosted material, and research preprints, with bibliographic checks performed on September 21, 2026. Searches included combinations of ``Hahn,'' ``separately holomorphic,'' ``Weierstrass,'' ``parameters,'' ``surcomplex,'' ``finite,'' and ``residue.'' Some searches were uninformative; absence of a matching hit is not evidence of absence from the literature. Accordingly, the results above are proposed, precisely scoped contributions, not certified priority claims.

\section{What has been resolved, and what remains outside the proof}\label{sec:limits}
\subsection{The part of the source program settled here}
Within the category defined in \cite{Source}, the following questions have affirmative answers under the hypotheses proved above. Separate Hahn coherence forces joint coherence. Parameter preparation and division can be carried out on one ordinary domain. Arbitrary coupled positive-support perturbations of coordinate powers define finite-free algebras with an exact rectangular basis. Their fibers consist of actual surcomplex roots with conserved formal local multiplicity. A multidimensional coefficientwise residue functional supplies a perfect duality across collisions.

These statements provide a concrete several-variable finite-mapping and finite-intersection theory. They are materially stronger than merely asserting that each ordinary coefficient satisfies a familiar theorem: the common support in the Hartogs theorem is not given in advance, and the coupled finite-algebra proof requires compatible quotient operations and uniform support control through matrix inversions.

\subsection{Hypotheses not to discard}
Positivity of the errors is essential to the proofs. For instance, changing the leading system to one with a common positive-dimensional zero set does not give a finite algebra of rank \(N\). If the coefficients are only fine-locally analytic with unrelated germs on different monads, the ordinary identity principle and the Baire-support proof do not apply. If a separate slice is known coherent only at ordinary values of the other variables, its behavior at arbitrary nonordinary tuples remains unprescribed.

Uniform ordinary domains are also part of the conclusion, not an empty formality. An infinite sequence of choices of smaller neighborhoods need not retain a usable ordinary neighborhood. The division operators in \cref{lem:ordinarydivision} and the induced matrix construction were arranged specifically to avoid that issue.

Finally, the positive-support hypothesis belongs to the coherent data, not merely to an informal assertion that an error is small at selected points. Theorems stated in terms of \(\Ip(V)\) carry a uniform global support condition. They do not automatically apply to arbitrary pointwise infinitesimal functions.

\subsection{Further directions, stated without a priority claim}
A natural next extension replaces the coordinate-power leading system by an arbitrary ordinary finite complete intersection, with a finite holomorphic basis over the parameter domain. To reproduce the present proof, one would need a fixed-domain ordinary division complex or compatible splitting that can be perturbed with positive support. Local algebra alone does not supply the required global support bookkeeping.

A second comparison problem is to identify the multidimensional localized functionals of \cref{sec:residue} with a chosen general Grothendieck-residue construction, including the Jacobian trace identity in that formulation. The present perfect-pairing theorem is independent of that comparison, and the coupled and one-variable examples provide exact normalization tests.

A third direction concerns coherent branch loci over several parameters. The finite multiplication matrices furnish characteristic polynomials and discriminants as coherent invariants, but a complete treatment of normalization, ramification strata, or resolution would require additional geometry. The finite-flat algebra theorem is not a substitute for those results.

None of these directions resolves the source manuscript's distinct global problem of sectorial or asymptotic re-expansion at arbitrarily large surreal scales. The coefficientwise finite-domain theory deliberately leaves that problem separate.

\section{Conclusion}
The main construction is a support-controlled deformation of a finite ordinary division problem. Positive Hahn supports replace an analytic convergence estimate, fixed ordinary division operators replace repeated shrinking, and induced multiplication matrices handle the coupling of several equations. The outcome is an exact finite-free algebra, an identification of its factors with actual infinitesimal local intersection algebras, and a residue pairing whose invertibility is visible already in its leading coefficient matrix.

The automatic-support theorem complements this algebraic construction: it explains when apparently unrelated one-variable coherent slices must come from one several-variable Hahn datum. Together these results supply a substantive, explicitly delimited extension of the supplied surcomplex framework. Their proofs and examples are intended for independent mathematical scrutiny, with the support conditions and limitations made part of the statements rather than left implicit.

\clearpage
\appendix
\section{Support and dependency audit}\label{app:support}
All supports in this appendix are set-sized. \(S\) denotes an input support; \(E\) a well-ordered positive error support; \(M=E^*\); and \(E_\eta\) the union of the positive supports of a finite displacement or parameter correction.

\begin{longtable}{@{}L{0.37\textwidth}L{0.23\textwidth}L{0.29\textwidth}@{}}
\toprule
Construction & Controlling support & Ordinary domain behavior\\
\midrule
\endhead
Evaluation or parameter pullback & \(S+E_\eta^*\) & Coefficients remain on the target pullback domain.\\[0.4em]
Positive operator inverse & \(S+E^*\) & Fixed input coefficient space is preserved.\\[0.4em]
Normalized preparation factors & \(E^*\) & One ordinary shrink before recursion; none during it.\\[0.4em]
Matrix quotient and remainder & \(S+M\) & Same parameter domain and fiber disk.\\[0.4em]
Induced multiplication matrices & \(M\) & Same remaining-variable polydisk.\\[0.4em]
Coupled normal form and quotient witnesses & \(S+M\) & Original product domain throughout.\\[0.4em]
Multidimensional residue & \(S+M\) & Holomorphic on the same parameter domain.\\[0.4em]
Residue Gram matrix and its inverse & \(M\) & No discriminant denominator or shrink.\\[0.4em]
Trace element and structure constants & \(M\) & Same parameter domain.\\[0.4em]
Individual root idempotents or weights & May require larger supports and negative exponents & Used in a fixed algebraically closed fiber workspace; no integral bound is asserted.\\
\bottomrule
\end{longtable}

At a fixed output exponent, every construction above that claims strong summability has only finitely many contributing support words. At every induction stage there are only finitely many matrix entries and basis monomials. These two separate finiteness statements justify regrouping all displayed sums. In particular, the support estimate for a global residue does not imply that its individual local summands have nonnegative support; \eqref{eq:fourpointsresidues} is an explicit counterexample to that inference.

\clearpage
\section{Verification artifacts and their limits}\label{app:verification}
The accompanying \texttt{verify\_examples.py} uses exact symbolic arithmetic to check:
\begin{enumerate}
\item the coupled multiplication matrices, their commutation and defining equations, and their characteristic polynomials;
\item the residue Gram matrix, dual basis, trace element, and root-weight identities on the four-element basis;
\item the displayed correction coefficients of the nonpolynomial sine system, modulo \(\zeta^3-1\), through the orders stated in the report;
\item the cubic residue and trace determinants, the repeated-root factorization, and the partial-fraction formula;
\item a low-order matrix-division example with a nonconstant analytic numerator and a genuinely matrix-valued error.
\end{enumerate}
The generated \texttt{verification\_report.txt} records the checks actually run and their symbolic truncation orders. These are finite tests of formulas. They do not verify transfinite recursions, arbitrary supports, the classical Hartogs input, or the identification of all local algebras. Those assertions rely on the proofs in the article.

The LaTeX file is standalone and contains its bibliography. Compile it with \texttt{latexmk -pdf surcomplex\_research.tex}, or run \texttt{pdflatex} repeatedly until references stabilize. No source image, external data file, or bibliography processor is needed to produce the PDF.

\section{Notation reference}
{\small
\begin{longtable}{@{}L{0.29\textwidth}L{0.62\textwidth}@{}}
\toprule
Symbol & Meaning\\
\midrule
\endhead
\(\SC=\No[i]\) & Surcomplex class field.\\[0.25em]
\(t^\gamma=\omega^{-\gamma}\) & Conway normal-form monomial.\\[0.25em]
\(K_\Gamma=\C((t^\Gamma))\) & Set-sized Hahn workspace.\\[0.25em]
\(U\sh\) & Finite surcomplex points with standard part in \(U\).\\[0.25em]
\(\HH(U)\) & Hahn-coherent data with ordinary holomorphic coefficient functions.\\[0.25em]
\(\Hp(U),\Ip(U)\) & Nonnegative-support ring and its positive-support ideal.\\[0.25em]
\(E^*\) & Positive-support additive monoid, including zero.\\[0.25em]
\(D_0,R_0\) & Ordinary quotient and remainder operators for a reference polynomial.\\[0.25em]
\(\boxset\) & Rectangular multi-index set \(0\leq\alpha_i<m_i\).\\[0.25em]
\(\mathcal A,\mathcal A_a,\mathcal A_{a,\xi}\) & Relative finite algebra, specialized algebra, and local root factor.\\[0.25em]
\(\mu_\xi(F_a)\) & Dimension of the formal local intersection algebra.\\[0.25em]
\(\Lambda_F\) & Coefficientwise multidimensional residue functional.\\[0.25em]
\(G,b^\alpha\) & Residue Gram matrix and residue-dual basis.\\[0.25em]
\(\mathfrak e_F\) & Element representing the trace functional under residue duality.\\
\bottomrule
\end{longtable}

}
\clearpage
\begin{thebibliography}{99}
\small
\addcontentsline{toc}{section}{References}

\bibitem{Source}
\emph{Surcomplex Analysis: Infinitesimal Calculus, Hahn-Coherent Holomorphy, and Contour Theory}.
Supplied research manuscript, file \texttt{surcomplex\_analysis(3).tex}, dated September 21, 2026.
In particular, the sections ``Hahn-coherent holomorphic functions,'' ``Weierstrass preparation and exact infinitesimal root lifting,'' ``Residues at moving poles, the argument principle, and Rouch\'e,'' and ``Further mathematical targets suggested by the results.''
This item is the user-supplied source, not an independently verified publication.

\bibitem{Alling}
Norman L. Alling,
\emph{Foundations of Analysis over Surreal Number Fields},
North-Holland Mathematics Studies \textbf{141}, Notas de Matem\'atica \textbf{117}, North-Holland, 1987.
\href{https://shop.elsevier.com/books/foundations-of-analysis-over-surreal-number-fields/alling/978-0-444-70226-5}{Publisher record}.

\bibitem{vdDE}
Lou van den Dries and Philip Ehrlich,
``Fields of surreal numbers and exponentiation,''
\emph{Fundamenta Mathematicae} \textbf{167} (2001), 173--188.
\href{https://doi.org/10.4064/fm167-2-3}{doi:10.4064/fm167-2-3}.

\bibitem{Poonen}
Bjorn Poonen,
``Maximally complete fields,''
\emph{L'Enseignement Math\'ematique} (2) \textbf{39} (1993), 87--106.
\href{https://math.mit.edu/~poonen/papers/amsval.pdf}{Author-hosted text}.

\bibitem{BM}
Alessandro Berarducci and Vincenzo Mantova,
``Transseries as germs of surreal functions,''
\emph{Transactions of the American Mathematical Society} \textbf{371} (2019), 3549--3592.
\href{https://doi.org/10.1090/tran/7428}{doi:10.1090/tran/7428};
\href{https://arxiv.org/abs/1703.01995}{arXiv:1703.01995}.

\bibitem{CL}
Raf Cluckers and Leonard Lipshitz,
``Fields with analytic structure,''
\emph{Journal of the European Mathematical Society} \textbf{13} (2011), no.~4, 1147--1223.
\href{https://doi.org/10.4171/JEMS/278}{doi:10.4171/JEMS/278}.

\bibitem{CLstrict}
Raf Cluckers and Leonard Lipshitz,
``Strictly convergent analytic structures,''
\emph{Journal of the European Mathematical Society} \textbf{19} (2017), no.~1, 107--149.
\href{https://doi.org/10.4171/JEMS/662}{doi:10.4171/JEMS/662}.

\bibitem{MS}
J\"orn M\"uller and Alexander Strohmaier,
``The theory of Hahn-meromorphic functions, a holomorphic Fredholm theorem, and its applications,''
\emph{Analysis \& PDE} \textbf{7} (2014), no.~3, 745--770.
\href{https://doi.org/10.2140/apde.2014.7.745}{doi:10.2140/apde.2014.7.745};
\href{https://arxiv.org/abs/1205.0236}{arXiv:1205.0236}.

\bibitem{Lebl}
Ji\v r\'i Lebl,
\emph{Tasty Bits of Several Complex Variables}, version 4.4, May 31, 2026.
Author's textbook, especially Theorems 4.3.1 and 6.2.5.
\href{https://www.jirka.org/scv/}{Author's book page};
\href{https://www.jirka.org/scv/scv.pdf}{PDF}.

\bibitem{BK}
Jacek Bochnak and Wojciech Kucharz,
``Global variants of Hartogs' theorem,''
\emph{Archiv der Mathematik} \textbf{113} (2019), 281--290.
\href{https://doi.org/10.1007/s00013-019-01340-7}{doi:10.1007/s00013-019-01340-7}.
The classical separate-holomorphy theorem recalled there is the background input in \cref{thm:hartogs}; the present proof does not use the article's stronger algebraic variants.

\end{thebibliography}
\end{document}
